%! Principles of Experimental Design — Unit 1, Lesson 1.6 — Exit Ticket
\documentclass[10pt]{article}
\usepackage{saar-article}
\usepackage{saar-boxes}
\newcommand{\TallMath}[1]{$\displaystyle #1\rule[-1.4em]{0pt}{3.2em}$}

\begin{document}

\pageheader{Unit 1, Lesson 1.6}{Exit Ticket}
% NAMESTRIP: no \namedateperiod here. The student writes their name once, on the
% cover this component is stapled behind. See references/conventions.md.

\vspace{0.15in}

\begin{notesbox}{Exit Ticket --- No Notes}
\small
Bayside Middle School tests a $10$-minute silent reading block. $80$
sixth-graders volunteer, and a random number generator assigns $40$ of them to
the reading block and the other $40$ to their usual homeroom, for one quarter.
Afterward, $28$ of the reading-block students improved their reading score, and
$16$ of the homeroom students did.

\begin{enumerate}[leftmargin=1.6em, itemsep=8pt, topsep=3pt]

  \item Nobody let a student choose.

        \textbf{(a)} Is this an experiment? \blank{2cm} \quad Who decided the
        groups? \blank{3.6cm}

        \textbf{(b)} Treatment: \blank{3.8cm} \hfill Control group:
        \blank{3.8cm}

  \item Compute both rates.

        \textbf{(a)} What percent of the reading-block students improved?

        \begin{work}
          \dfrac{28}{40} &= 0.70 \\
                         &= 70\%
        \end{work}

        \textbf{(b)} What percent of the homeroom students improved?

        \begin{work}
          \dfrac{16}{40} &= 0.40 \\
                         &= 40\%
        \end{work}

        \textbf{(c)} The reading-block group is higher by \blank{2cm}
        percentage points.

  \item A teacher suggests a change: \emph{``Next year, just give the reading
        block to the students who want it.''} Name the principle that change
        breaks, and say what could go wrong because of it.
        \writelines{2}

\end{enumerate}
\end{notesbox}

\end{document}
